% Lösungen Wahrscheinlichkeit (zusammenbau v0.4, Heft)
\einheitenkopf{Wahrscheinlichkeit}

\erg{6}{a) 4 Paare: (grau; gelb), (grau; weiß), (grün; gelb), (grün; weiß) \quad b) $3 \cdot 2 = 6$ Kombinationen \quad c) 6 Ergebnisse: (Kopf; 1), (Kopf; 2), (Kopf; 3), (Zahl; 1), (Zahl; 2), (Zahl; 3)}
\erg{7}{a) 6 Zahlen: 147, 174, 417, 471, 714, 741; gerade sind 174 und 714, Anteil $\frac{2}{6} = \frac{1}{3}$ \quad b) 751 \quad c) 4 Ergebnisse: (rot; rot), (rot; blau), (blau; rot), (blau; blau) \quad d) 9 Zahlen: 44, 45, 46, 54, 55, 56, 64, 65, 66 \quad e) 6 Paare: (1; 4), (2; 4), (3; 4), (4; 4), (5; 4), (6; 4) \quad f) 9 \quad g) $3 \cdot 2 \cdot 1 = 6$ Codes (248, 284, 428, 482, 824, 842) \quad h) $3 \cdot 2 \cdot 1 = 6$ Codes}
\erg{8}{a) möglich: 35 Lose (30 Nieten und 5 Gewinne), günstig: 5 Gewinne \quad b) $\frac{1}{6}$ \quad c) $\frac{2}{8} = \frac{1}{4}$}
\erg{9}{a) $\frac{9}{50} = 0{,}18 = 18\,\%$ \quad b) $\frac{5}{8} = 0{,}625 = 62{,}5\,\%$ \quad c) $\frac{1}{8} = 0{,}125 = 12{,}5\,\%$ \quad d) $\frac{15}{45 + 15} = \frac{15}{60} = \frac{1}{4}$ \quad e) Ja: $\frac{1}{350}$, denn $600 - 251 + 1 = 350$ Lose, eines gewinnt. \quad f) $\frac{4}{60} = \frac{1}{15}$ \quad g) $\frac{6}{8} = \frac{3}{4}$ \quad h) $0{,}75 \cdot 4 = 3$ Felder}
\erg{10}{a) 8 Felder insgesamt, also 2 weiße Felder \quad b) Nein: $\frac{3}{16}$, denn es sind nur noch 16 Muffins da, davon 3 mit Chili. \quad c) Ja: $\frac{4}{50}$, denn Lose mit Endziffer 5 sind 205, 215, …, 695, also 50 Lose; davon enden auf 45: 245, 345, 445, 545, 645 – ohne 445 bleiben 4. \quad d) Nein: Lose mit Endziffer 2 sind 202, 212, …, 792, also 60 Lose; davon enden auf 72: 272, 372, 472, 572, 672, 772 – ohne 572 bleiben 5, also $\frac{5}{60} = \frac{1}{12}$.}

10a)

\kreisdiagramm{rot/1,rot/1,rot/1,rot/1,rot/1,rot/1,weiß/1,weiß/1}

\erg{11}{a) Mond}
\erg{12}{a) mittlere Dose}
\erg{13}{a) ohne Zurücklegen \quad b) an jedem Ast $\frac{1}{2}$ \quad c) erste Stufe daneben $\frac{1}{5}$; nach Treffer: Treffer $\frac{4}{5}$; nach daneben: daneben $\frac{1}{5}$}

13b)

\baumzwei{Kopf/\frac{1}{2},Zahl/\frac{1}{2}}{Kopf/\frac{1}{2},Zahl/\frac{1}{2}}{Kopf/\frac{1}{2},Zahl/\frac{1}{2}}


13c)

\baumzwei{Treffer/\frac{4}{5},daneben/\frac{1}{5}}{Treffer/\frac{4}{5},daneben/\frac{1}{5}}{Treffer/\frac{4}{5},daneben/\frac{1}{5}}

\erg{14}{a) grüner Würfel ungerade $\frac{3}{6}$; gelber Würfel an beiden Punkten gerade $\frac{5}{6}$, ungerade $\frac{1}{6}$ \quad b) an jedem Punkt grün $\frac{3}{5}$, gelb $\frac{2}{5}$ – das Glücksrad ändert sich nicht \quad c) $\frac{1}{3} \cdot \frac{1}{3} = \frac{1}{9}$ – jede Woche liegen wieder alle drei Zettel im Topf. \quad d) $\frac{3}{4} \cdot \frac{2}{4} + \frac{1}{4} \cdot \frac{1}{4} = \frac{6}{16} + \frac{1}{16} = \frac{7}{16}$ \quad e) $\frac{25}{64} + \frac{4}{64} + \frac{1}{64} = \frac{30}{64} = \frac{15}{32}$ \quad f) $1 - \frac{3}{6} \cdot \frac{5}{6} = 1 - \frac{15}{36} = \frac{21}{36} = \frac{7}{12}$}

14a)

\baumzwei{gerade/\frac{3}{6},ungerade/\frac{3}{6}}{gerade/\frac{5}{6},ungerade/\frac{1}{6}}{gerade/\frac{5}{6},ungerade/\frac{1}{6}}


14b)

\baumzwei{grün/\frac{3}{5},gelb/\frac{2}{5}}{grün/\frac{3}{5},gelb/\frac{2}{5}}{grün/\frac{3}{5},gelb/\frac{2}{5}}

\erg{15}{a) Nein: 66 hat $\frac{1}{4} \cdot \frac{2}{4} = \frac{2}{16}$, 77 hat $\frac{1}{4} \cdot \frac{1}{4} = \frac{1}{16}$. \quad b) Nein: Zwei gleiche Augenzahlen haben $\frac{6}{36} = \frac{1}{6}$, eine 1 im ersten Wurf hat auch $\frac{1}{6}$ – beide sind gleich wahrscheinlich. \quad c) z. B. links 5, 5, 5, 6 und rechts 7, 7, 8, 8: $\frac{3}{4} \cdot \frac{2}{4} = \frac{6}{16} = \frac{3}{8}$ \quad d) z. B. eine 3 durch eine 2 ersetzen: $\frac{2}{6} \cdot \frac{2}{6} = \frac{4}{36} = \frac{1}{9}$ \quad e) an jedem Ast Stern $\frac{1}{4}$, leer $\frac{3}{4}$; $3 \cdot \frac{1}{4} \cdot \frac{1}{4} \cdot \frac{3}{4} + \frac{1}{4} \cdot \frac{1}{4} \cdot \frac{1}{4} = \frac{9}{64} + \frac{1}{64} = \frac{10}{64} = \frac{5}{32}$ \quad f) an jedem Ast größer als 4 $\frac{1}{3}$, höchstens 4 $\frac{2}{3}$; $3 \cdot \frac{1}{3} \cdot \frac{1}{3} \cdot \frac{2}{3} + \frac{1}{3} \cdot \frac{1}{3} \cdot \frac{1}{3} = \frac{6}{27} + \frac{1}{27} = \frac{7}{27}$}

15c)

\kreisdiagramm{5/1,5/1,5/1,6/1} \kreisdiagramm{7/1,7/1,8/1,8/1}


15e)

\baumdreigleich{Stern/\frac{1}{4},leer/\frac{3}{4}}


15f)

\baumdreigleich{größer als 4/\frac{1}{3},höchstens 4/\frac{2}{3}}

\erg{16}{a) 8 Kekse, davon 3 mit Schoko \quad b) $\frac{3}{8} \cdot \frac{2}{7} = \frac{6}{56} = \frac{3}{28}$ \quad c) erste Stufe grün $\frac{3}{7}$; nach rot: rot $\frac{3}{6}$; nach grün: rot $\frac{4}{6}$, grün $\frac{2}{6}$}

16c)

\baumzwei{rot/\frac{4}{7},grün/\frac{3}{7}}{rot/\frac{3}{6},grün/\frac{3}{6}}{rot/\frac{4}{6},grün/\frac{2}{6}}

\erg{17}{a) $\frac{6}{7} \cdot \frac{1}{6} = \frac{1}{7}$ \quad b) erste Stufe ohne Nuss $\frac{10}{15}$; nach ohne, Nuss: Nuss $\frac{4}{13}$ \quad c) $\frac{1}{9} + \frac{8}{9} \cdot \frac{1}{8} + \frac{8}{9} \cdot \frac{7}{8} \cdot \frac{1}{7} = \frac{3}{9} = \frac{1}{3}$ \quad d) Drei rote: $\frac{5}{15} \cdot \frac{4}{14} \cdot \frac{3}{13} = \frac{60}{2\,730}$, drei grüne: $\frac{3}{15} \cdot \frac{2}{14} \cdot \frac{1}{13} = \frac{6}{2\,730}$ – zehnmal so wahrscheinlich, nicht doppelt. \quad e) (Vanille; Pflaume), (Pflaume; Vanille), (Vanille; Vanille); $\frac{15}{18} \cdot \frac{14}{17} = \frac{210}{306} = \frac{35}{51} \approx 0{,}686$; Ereignis: beide Berliner haben die gleiche Füllung. \quad f) (Marzipan; Nougat), (Nougat; Marzipan), (Marzipan; Marzipan); $\frac{13}{15} \cdot \frac{12}{14} = \frac{156}{210} = \frac{26}{35} \approx 0{,}743$; Ereignis: genau eine Praline mit Marzipan.}

17b)

\baumdrei{Nuss/\frac{5}{15},ohne/\frac{10}{15}}{Nuss/\frac{4}{14},ohne/\frac{10}{14}}{Nuss/\frac{5}{14},ohne/\frac{9}{14}}{Nuss/\frac{3}{13},ohne/\frac{10}{13}}{Nuss/\frac{4}{13},ohne/\frac{9}{13}}{Nuss/\frac{4}{13},ohne/\frac{9}{13}}{Nuss/\frac{5}{13},ohne/\frac{8}{13}}

